\documentclass[12pt, a4paper]{article}
 
\usepackage[margin=2.5cm]{geometry}
\usepackage{amsmath, amssymb}
\usepackage{graphicx}
\usepackage{hyperref}
\usepackage{booktabs}
\usepackage{caption}
\usepackage{subcaption}
\usepackage{setspace}
\usepackage{parskip}
\usepackage{natbib}
\usepackage{microtype}
\usepackage{xcolor}
 
\hypersetup{
    colorlinks=true,
    linkcolor=blue!60!black,
    citecolor=blue!60!black,
    urlcolor=blue!60!black
}
 
\onehalfspacing
 
\title{\textbf{A Weak Lensing Forward Model for Flat $\Lambda$CDM}\\
\large From the Matter Power Spectrum to Shear Correlation Functions}
 
\author{Ashok Kumar Ramu\\
\small Department of Physics and Astronomy, University of Waterloo}
 
\date{March 2026}
 
\begin{document}
 
\maketitle
 
\begin{abstract}
I build a weak gravitational lensing forward model from scratch in Python using
only \texttt{numpy} and \texttt{scipy}, with no external cosmology libraries.
Starting from the FLRW background, I implement the Eisenstein \& Hu (1998)
transfer function, the Heath (1977) growth factor, and the lensing convergence
kernel to compute the shear power spectrum $C_\ell$ via the Limber approximation.
The shear correlation functions $\xi_\pm(\theta)$ are recovered by Hankel
transform. I reproduce the observed $S_8$ tension between Planck CMB constraints
and weak lensing surveys (KiDS-1000, DES Y3) at the $\sim3.8\sigma$ level, and
characterize the scale-dependent sensitivity of $\xi_+$ to $\Omega_m$ and
$\sigma_8$. The main result is that $\sigma_8$ shifts the signal by a uniform
$\sim$21\% per 10\% parameter change across all angular scales, while $\Omega_m$
produces a scale-dependent response that changes sign near $\theta \sim 200$
arcmin. This is directly relevant to the measurement strategy of wide-field
lensing surveys such as UNIONS and Euclid.
\end{abstract}

\newpage
\tableofcontents
\newpage
 
% ============================================================
\section{Motivation}
% ============================================================
 
Weak gravitational lensing is one of the cleanest probes of the matter
distribution in the Universe. Background galaxies appear slightly elliptical
because their light follows geodesics deflected by the foreground metric
perturbations sourced by large-scale structure. Measuring these coherent
distortions across the sky gives a direct statistical handle on the amplitude
and shape of the matter power spectrum, and therefore on $\Omega_m$ and
$\sigma_8$.
 
The reason this is interesting right now is the $S_8$ tension.
The combination $S_8 \equiv \sigma_8\sqrt{\Omega_m/0.3}$ is what weak lensing
surveys constrain most tightly, and current measurements sit $\sim3$--$4\sigma$
below the value inferred from the Planck CMB. The Kilo-Degree Survey
\citep{Asgari2021} finds $S_8 = 0.766 \pm 0.017$ and the Dark Energy Survey
Year 3 analysis \citep{Amon2022} gives $S_8 = 0.776 \pm 0.017$, both
systematically lower than the Planck 2018 value of $S_8 = 0.831 \pm 0.006$
\citep{Aghanim2020}. Whether this is a systematic effect or new physics is
an open question.
 
The UNIONS survey \citep[Ultraviolet Near-Infrared Northern Survey;][]{UNIONS},
currently being carried out with the Canada-France-Hawaii Telescope, will provide
an independent northern-hemisphere lensing measurement and is expected to
contribute to resolving this tension. The motivation for this project is to
understand, from first principles, what UNIONS is actually measuring before
attempting any parameter inference on top.
 
The goal here is not to produce a production-level pipeline. It is to build
an honest forward model---one where every line of physics is visible and
testable---and use it to understand which angular scales carry cosmological
information and how.
 
% ============================================================
\section{Theory}
% ============================================================
 
\subsection{FLRW Background}
 
I work throughout in flat $\Lambda$CDM. The dimensionless Hubble parameter is
\begin{equation}
    E(z) \equiv \frac{H(z)}{H_0} = \sqrt{\Omega_m(1+z)^3 + \Omega_\Lambda},
\end{equation}
where $\Omega_\Lambda = 1 - \Omega_m$ enforces flatness. The comoving distance
to redshift $z$ is
\begin{equation}
    \chi(z) = \frac{c}{H_0} \int_0^z \frac{dz'}{E(z')},
\end{equation}
which I evaluate numerically using a 500-point trapezoidal rule at each $z$.
As a sanity check, $\chi(z=1100) \approx 14{,}000\,\mathrm{Mpc}$, consistent
with the comoving distance to the CMB surface of last scattering.
 
\subsection{Growth Factor}
 
Structure growth in the linear regime is captured by the growth factor $D(z)$,
normalized so that $D(0) = 1$. I use the integral form from \citet{Heath1977},
\begin{equation}
    D(z) \propto H(z)\int_z^\infty \frac{1+z'}{H^3(z')}\,dz',
\end{equation}
evaluated with \texttt{scipy.integrate.quad} and normalized to unity at $z=0$.
 
\subsection{Matter Power Spectrum}
 
The linear matter power spectrum at $z=0$ is
\begin{equation}
    P(k) = A_s\, k^{n_s}\, T^2(k),
\end{equation}
where $n_s = 0.9649$ is the spectral index \citep{Aghanim2020} and $T(k)$ is
the transfer function. I use the Eisenstein \& Hu (1998) fitting formula
\citep{EisensteinHu1998}, which handles baryon acoustic features and is accurate
to $\sim$5\% for standard $\Lambda$CDM. The simpler BBKS approximation ignores
baryons entirely and would introduce a visible error at the percent level.
 
The normalization $A_s$ is set by requiring
\begin{equation}
    \sigma_8^2 = \frac{1}{2\pi^2}\int_0^\infty k^2\,P(k)\,|\tilde{W}(kR)|^2\,dk,
    \quad R = 8\,h^{-1}\,\mathrm{Mpc},
\end{equation}
where $\tilde{W}(x) = 3(\sin x - x\cos x)/x^3$ is the Fourier transform of a
real-space top-hat window. Varying $\sigma_8$ at fixed $\Omega_m$ rescales
$P(k)$ by a constant factor, which is why it affects $\xi_\pm$ uniformly across
all angular scales (Section~\ref{sec:sensitivity}).
 
At redshift $z$, the power spectrum is suppressed by the growth factor:
$P(k,z) = P(k,0)\,D^2(z)$.
 
\subsection{Weak Lensing Kernel}
 
The lensing convergence $\kappa$ is a weighted projection of the matter overdensity
along the line of sight:
\begin{equation}
    \kappa(\hat{n}) = \int_0^{\chi_H} W(\chi)\,\delta(\chi\hat{n},\,z(\chi))\,d\chi,
\end{equation}
where the lensing kernel is
\begin{equation}
    W(\chi) = \frac{3}{2}\left(\frac{H_0}{c}\right)^2\Omega_m
    \frac{\chi}{a(\chi)}\int_\chi^{\chi_H}
    n(\chi')\frac{\chi'-\chi}{\chi'}\,d\chi'.
    \label{eq:kernel}
\end{equation}
The prefactor $(3/2)(H_0/c)^2\Omega_m$ comes from combining the Poisson equation
in comoving coordinates with the GR light deflection formula; most review papers
quote only the final Limber result without showing this derivation.
 
The source galaxy redshift distribution $n(z)$ is modelled as
\begin{equation}
    n(z) \propto \left(\frac{z}{z_*}\right)^\alpha
    \exp\!\left[-\left(\frac{z}{z_*}\right)^\beta\right],
\end{equation}
with $z_* = 0.5$, $\alpha = 1.5$, $\beta = 2.0$, normalized to unit integral.
This is a plausible approximation for a ground-based survey at SDSS/UNIONS depth;
I do not use the actual UNIONS $n(z)$ here.
 
\subsection{Shear Power Spectrum and Correlation Functions}
 
The shear $E$-mode power spectrum follows from the Limber approximation
\citep{Limber1953,Loverde2008},
\begin{equation}
    C_\ell = \int_0^{\chi_H}
    \frac{W^2(\chi)}{\chi^2}\,P\!\left(\frac{\ell+1/2}{\chi},\,z(\chi)\right)d\chi,
    \label{eq:limber}
\end{equation}
which is accurate for $\ell \gg 1$, i.e.\ throughout the range I use
($\ell = 2$--$2000$). I evaluate this on a geometrically-spaced grid of 80
multipoles.
 
The shear correlation functions are related to $C_\ell$ by Hankel transforms:
\begin{align}
    \xi_+(\theta) &= \frac{1}{2\pi}\int_0^\infty \ell\,C_\ell\,J_0(\ell\theta)\,d\ell,
    \label{eq:xip}\\
    \xi_-(\theta) &= \frac{1}{2\pi}\int_0^\infty \ell\,C_\ell\,J_4(\ell\theta)\,d\ell,
    \label{eq:xim}
\end{align}
where $J_0$ and $J_4$ are Bessel functions of the first kind. Both $J_0$ and $J_4$ appear because $\xi_+$ and $\xi_-$ are two distinct 
components of the shear two-point function, related to the $E$-mode power 
spectrum by separate Hankel transforms of different order. The orders 0 and 
4 follow from the spin-2 nature of the shear field 
\citep[eq.~2.69]{Bartelmann2001}; they are not interchangeable. These integrals are evaluated on a dense
grid of $8000$ multipoles using \texttt{scipy.special.jv} and
\texttt{numpy.trapezoid}.
 
% ============================================================
\section{Results}
% ============================================================
 
I use Planck 2018 best-fit parameters as a fiducial throughout:
$\Omega_m = 0.315$, $\sigma_8 = 0.811$, $\Omega_b = 0.0486$,
$H_0 = 67.3\,\mathrm{km\,s^{-1}\,Mpc^{-1}}$, $n_s = 0.9649$
\citep{Aghanim2020}.
 
\subsection{Matter Power Spectrum}
 
Figure~\ref{fig:pk} shows $P(k)$ at $z=0$ for a range of $\Omega_m$ (left) and
$\sigma_8$ (right). Increasing $\Omega_m$ shifts the turnover to larger $k$
because matter-radiation equality occurs at a higher redshift, compressing the
horizon scale. Increasing $\sigma_8$ uniformly rescales the amplitude with no
change to the shape, as expected from equation~(\ref{eq:limber}).
 
\begin{figure}[h]
    \centering
    \includegraphics[width=\linewidth]{fig1_power_spectrum.png}
    \caption{Linear matter power spectrum $P(k)$ at $z=0$ computed using the
    Eisenstein \& Hu (1998) transfer function. \textit{Left:} Varying $\Omega_m$
    with $\sigma_8 = 0.811$ fixed. Increasing $\Omega_m$ shifts the turnover
    scale to larger $k$ and suppresses large-scale power. \textit{Right:} Varying
    $\sigma_8$ with $\Omega_m = 0.315$ fixed. $\sigma_8$ is a pure normalization
    with no effect on the shape. The Planck 2018 fiducial is shown as a solid
    orange curve in both panels.}
    \label{fig:pk}
\end{figure}
 
\subsection{Lensing Kernel}
 
Figure~\ref{fig:kernel} shows the lensing kernel $W(\chi)$ from
equation~(\ref{eq:kernel}) for three values of $\Omega_m$. Two effects are
visible. First, the amplitude scales roughly as $\Omega_m^2$: one factor comes
from the prefactor in equation~(\ref{eq:kernel}), and a second from the larger
density contrast $\delta$ in a more matter-dominated universe. Second, the
peak shifts to lower $\chi$ for higher $\Omega_m$ because the comoving
distance--redshift relation is compressed.
 
\begin{figure}[h]
    \centering
    \includegraphics[width=0.78\linewidth]{fig2_lensing_kernel.png}
    \caption{Weak lensing kernel $W(\chi)$ for three values of $\Omega_m$,
    with the corresponding redshift scale shown on the upper axis. The kernel
    peaks near $\chi \sim 1000\,\mathrm{Mpc}$ ($z \sim 0.3$) for the Planck
    fiducial, reflecting the geometric weighting between observer, lens, and
    source planes.}
    \label{fig:kernel}
\end{figure}
 
\subsection{Shear Correlation Functions}
 
Figures~\ref{fig:xi_om} and~\ref{fig:xi_s8} show $\xi_+$ and $|\xi_-|$ at
eleven angular scales spanning $1$--$300$ arcmin. Both correlation functions
decrease monotonically with $\theta$ in the range plotted; the dip visible in
$|\xi_-|$ near $\theta \sim 2$--$5$ arcmin reflects a sign change in $\xi_-$
at those scales.
 
The curves separate cleanly for different parameter values in both figures,
confirming that the forward model is sensitive enough to distinguish $\Delta\Omega_m
\sim 0.05$ and $\Delta\sigma_8 \sim 0.1$ over this angular range. This is the
signal that UNIONS measures from galaxy ellipticities using TreeCorr
\citep{Jarvis2004}.
 
\begin{figure}[h]
    \centering
    \includegraphics[width=\linewidth]{fig3_xi_vary_omega_m.png}
    \caption{Shear correlation functions $\xi_+(\theta)$ (left) and
    $|\xi_-(\theta)|$ (right) for five values of $\Omega_m$ with
    $\sigma_8 = 0.811$ fixed.}
    \label{fig:xi_om}
\end{figure}
 
\begin{figure}[h]
    \centering
    \includegraphics[width=\linewidth]{fig4_xi_vary_sigma8.png}
    \caption{Same as Figure~\ref{fig:xi_om} but varying $\sigma_8$ with
    $\Omega_m = 0.315$ fixed. The separation between curves is more uniform
    across scales than for $\Omega_m$.}
    \label{fig:xi_s8}
\end{figure}
 
\subsection{The $S_8$ Tension}
 
Figure~\ref{fig:s8} shows the $(\Omega_m, \sigma_8)$ plane with contours of
constant $S_8 = \sigma_8\sqrt{\Omega_m/0.3}$ overlaid. The Planck 2018 best-fit
sits at $S_8 = 0.831 \pm 0.006$, while KiDS-1000 \citep{Asgari2021} and
DES Y3 \citep{Amon2022} both cluster near $S_8 \approx 0.77$. The tension
is $\sim3.8\sigma$ when comparing Planck to KiDS-1000.
 
The contour structure in Figure~\ref{fig:s8} is directly related to
the degeneracy visible in the correlation functions: at small angular scales
where both $\Omega_m$ and $\sigma_8$ produce a similar $\sim$21\% response
per 10\% parameter shift, lensing essentially constrains only the combination
$S_8$ and cannot separate the two parameters individually.
 
\begin{figure}[h]
    \centering
    \includegraphics[width=0.78\linewidth]{fig5_s8_tension.png}
    \caption{The $S_8$ tension in the $(\Omega_m, \sigma_8)$ plane. Dashed
    grey contours show lines of constant $S_8$; the colour map encodes the
    $S_8$ value. Planck 2018 (red star) sits $\sim3.8\sigma$ above the KiDS-1000
    (blue diamond) and DES Y3 (green triangle) measurements along the $S_8$
    degeneracy direction. Approximate $1\sigma$ ellipses are shown for each
    experiment.}
    \label{fig:s8}
\end{figure}
 
\subsection{Parameter Sensitivity}
\label{sec:sensitivity}
 
Figure~\ref{fig:sensitivity} quantifies the fractional change in $\xi_+$ when
$\Omega_m$ or $\sigma_8$ is shifted by 10\% from the Planck fiducial. Two
results stand out.
 
First, $\sigma_8$ produces a scale-independent response of $\pm$21\% across
all angular scales. This follows analytically from $\xi_+ \propto \sigma_8^2$:
a 10\% shift in $\sigma_8$ gives a $1.1^2 - 1 \approx 21\%$ change in
$\xi_+$ at every $\theta$.
 
Second, the $\Omega_m$ response is strongly scale-dependent. At small angles
($\theta \lesssim 10$ arcmin) a 10\% increase in $\Omega_m$ raises $\xi_+$ by
$\sim$21\%, similar to $\sigma_8$. This is the origin of the $S_8$ degeneracy:
at small scales the two parameters are nearly interchangeable. At large angles
($\theta \gtrsim 100$ arcmin) the $\Omega_m$ response drops toward zero and
changes sign near $\theta \sim 200$ arcmin. This sign change is a real feature
of the Hankel transform: as $\Omega_m$ increases, the power spectrum turnover
shifts and the large-angle correlation function actually decreases.
 
The practical implication is that angular scales above $\sim$100 arcmin carry
information about $\Omega_m$ that is \emph{independent} of the $S_8$ degeneracy
direction. Wide-field surveys like UNIONS, which cover thousands of square degrees,
have access to these large angular separations and can therefore place tighter
individual constraints on $\Omega_m$ than a narrow deep survey with the same
number of galaxies.
 
\begin{figure}[h]
    \centering
    \includegraphics[width=0.85\linewidth]{fig6_sensitivity.png}
    \caption{Fractional change in $\xi_+(\theta)$ for $\pm$10\% shifts in
    $\Omega_m$ (red) and $\sigma_8$ (blue), relative to the Planck fiducial.
    The $\sigma_8$ response is flat at $\pm$21\% by construction. The $\Omega_m$
    response is scale-dependent, peaking at small angles and changing sign near
    $\theta \sim 200$ arcmin. At small scales both parameters produce the same
    fractional change, which is the geometric origin of the $S_8$ degeneracy.}
    \label{fig:sensitivity}
\end{figure}
 
% ============================================================
\section{Discussion and Next Steps}
% ============================================================
 
The forward model reproduces known results: the $S_8$ tension at the correct
significance, the expected $\sigma_8^2$ scaling of the lensing signal, and the
qualitative shape of $\xi_\pm$ consistent with published KiDS and DES
measurements.

There are several approximations worth noting. The most significant is the use 
of the linear matter power spectrum throughout. At angular scales 
$\theta \lesssim 10$ arcmin, nonlinear structure growth substantially enhances 
the lensing signal relative to the linear prediction, and the two can differ by 
factors of several at arcminute scales. The sensitivity results in 
Figure~\ref{fig:sensitivity} should therefore be interpreted as valid only at 
$\theta \gtrsim 10$ arcmin. Implementing \texttt{HaloFit} \citep{Smith2003} or 
the \texttt{EuclidEmulator} \citep{EuclidEmulator2021} at small scales is the 
most important numerical extension of this work, and would be required before 
attempting any quantitative parameter inference against real data. The $n(z)$ used here is a smooth
analytic form; the real UNIONS $n(z)$ will have structure from photometric
redshift errors and tomographic binning. The Limber approximation in
equation~(\ref{eq:limber}) introduces errors of order $(1/\ell)^2$ at the
lowest multipoles but is sufficient for $\ell \gtrsim 10$. Nonlinear structure
growth, which becomes important below $\sim$10 arcmin, is not included; a
production pipeline would use \texttt{HaloFit} \citep{Smith2003} or the
\texttt{EuclidEmulator} \citep{EuclidEmulator2021} at those scales.
 
The natural next step is a Fisher matrix forecast for a UNIONS-like survey
geometry. Given the sensitivity derivatives already computed here, the Fisher
matrix
\begin{equation}
    F_{ij} = \sum_{\theta,\theta'}
    \frac{\partial\xi_+(\theta)}{\partial p_i}\,
    \text{Cov}^{-1}(\theta,\theta')\,
    \frac{\partial\xi_+(\theta')}{\partial p_j}
\end{equation}
would predict the expected $1\sigma$ constraints on $\Omega_m$ and $\sigma_8$
as a function of survey area and galaxy number density. The covariance matrix
for a survey of area $A$ and shape noise $\sigma_e^2/n_\mathrm{gal}$ is
well-approximated analytically \citep{Joachimi2008}.
 
A longer-term extension would be to train a neural network on a grid of
simulated $\xi_\pm$ vectors and use it for direct parameter inference on real
SDSS or UNIONS data. The sensitivity analysis in Section~\ref{sec:sensitivity}
suggests this is feasible: the $\xi_\pm$ vectors at 11 angular scales carry
sufficient information to distinguish cosmologies at the few-percent level in
both parameters.
 
% ============================================================
\section*{Code}
% ============================================================
 
The full notebook, including all figures, is available at
\url{https://github.com/Ashok1103/Physics-Research-Portfolio.git}.
All computations use only \texttt{numpy}, \texttt{scipy}, and
\texttt{matplotlib}. No external cosmology libraries are required.
 
% ============================================================
\begin{thebibliography}{99}

\bibitem[Aghanim et al.(2020)]{Aghanim2020}
Aghanim N., et al., 2020, A\&A, 641, A6

\bibitem[Asgari et al.(2021)]{Asgari2021}
Asgari M., et al., 2021, A\&A, 645, A104

\bibitem[Amon et al.(2022)]{Amon2022}
Amon A., et al., 2022, Phys. Rev. D, 105, 023514

\bibitem[Eisenstein \& Hu(1998)]{EisensteinHu1998}
Eisenstein D.~J., Hu W., 1998, ApJ, 496, 605

\bibitem[Heath(1977)]{Heath1977}
Heath D.~J., 1977, MNRAS, 179, 351

\bibitem[Bartelmann \& Schneider(2001)]{Bartelmann2001}
Bartelmann M., Schneider P., 2001, Phys. Rep., 340, 291

\bibitem[Limber(1953)]{Limber1953}
Limber D.~N., 1953, ApJ, 117, 134

\bibitem[LoVerde \& Afshordi(2008)]{Loverde2008}
LoVerde M., Afshordi N., 2008, Phys. Rev. D, 78, 123506

\bibitem[Jarvis et al.(2004)]{Jarvis2004}
Jarvis M., Bernstein G., Jain B., 2004, MNRAS, 352, 338

\bibitem[Smith et al.(2003)]{Smith2003}
Smith R.~E., et al., 2003, MNRAS, 341, 1311

\bibitem[Euclid Collaboration et al.(2021)]{EuclidEmulator2021}
Euclid Collaboration, Knabenhans M., et al., 2021, MNRAS, 505, 2840

\bibitem[Joachimi \& Schneider(2008)]{Joachimi2008}
Joachimi B., Schneider P., 2008, A\&A, 488, 829

\bibitem[Hudson et al.(2024)]{UNIONS}
Hudson M.~J., et al., 2024, UNIONS: Ultraviolet Near-Infrared Northern Survey,
\url{https://www.skysurvey.cc}

\end{thebibliography}

\end{document}
% ============================================================
 
\end{document}